\documentclass[11pt]{article}

\usepackage[T1]{fontenc}
\usepackage{lmodern}
\usepackage[a4paper,margin=30mm]{geometry}
\usepackage{amsmath,amssymb,amsthm}
\usepackage{enumitem}
\usepackage[expansion=false]{microtype}
\usepackage[hidelinks]{hyperref}
\hypersetup{
  pdftitle={Tuza's conjecture for line and total graphs},
  pdfauthor={Samuel Mass\'{e}},
  pdfsubject={Tuza's inequality for line graphs and total graphs},
  pdfkeywords={triangle packing, triangle transversal, line graph, total graph}
}

\newtheorem{theorem}{Theorem}[section]
\newtheorem{lemma}[theorem]{Lemma}
\newtheorem{corollary}[theorem]{Corollary}
\newtheorem{proposition}[theorem]{Proposition}
\theoremstyle{remark}
\newtheorem{remark}[theorem]{Remark}

\newcommand{\tri}{\triangle}
\newcommand{\Tot}{\operatorname{Tot}}
\newcommand{\dd}{\mathbin{\dot\cup}}

\title{Tuza's conjecture for line and total graphs}
\author{Samuel Mass\'{e}}
\date{11 September 2026}

\begin{document}

\maketitle

\begin{abstract}
For a graph $G$, let $\nu_\tri(G)$ be the maximum number of pairwise
edge-disjoint triangles and let $\tau_\tri(G)$ be the minimum number of edges
meeting every triangle.  We prove $\tau_\tri(G)\leq 2\nu_\tri(G)$ when $G$ is
the line graph or the total graph of any finite simple graph.  The proofs use
edge-disjoint incidence cliques and a two-coloring of root edges obtained
from Euler tours.  A local packing adjustment handles the exceptional
Euler-tour seam for line graphs; adjoining each original vertex to the
minority color of its incidence clique absorbs the seam for total graphs.
The factor two is sharp over line graphs.  A cover constructed from distinct
incident-edge representatives gives, for every cubic triangle-free graph $F$,
the exact identity
$\tau_\tri(\Tot(F))=\nu_\tri(\Tot(F))=5|V(F)|/2$.
The proof constructions presented here follow the accompanying Lean
formalization, including its explicit finite clique-packing certificates.
\end{abstract}

\section{Introduction}

A \emph{triangle packing} in a graph $G$ is a collection of pairwise
edge-disjoint triangles, and a \emph{triangle transversal} is a set of edges
meeting every triangle.  Their maximum and minimum cardinalities are denoted
by $\nu_\tri(G)$ and $\tau_\tri(G)$, respectively.  Tuza conjectured that
\begin{equation}\label{eq:tuza}
  \tau_\tri(G)\leq 2\nu_\tri(G)
\end{equation}
for every finite graph $G$ \cite{tuza1990}.

We establish \eqref{eq:tuza} for two graph operators built from the incidence
structure of a root graph.  The \emph{line graph} $L(H)$ has vertex set
$E(H)$, with two vertices adjacent when the corresponding root edges meet.
The \emph{total graph} $\Tot(F)$ has vertex set $V(F)\dd E(F)$; two elements
are adjacent when they are adjacent or incident in $F$.

Our main results are the following.

\begin{theorem}\label{thm:line}
If $H$ is a finite simple graph, then
\[
  \tau_\tri(L(H))\leq 2\nu_\tri(L(H)).
\]
The factor two cannot be improved over the class of line graphs.
\end{theorem}

\begin{theorem}\label{thm:total}
If $F$ is a finite simple graph, then
\[
  \tau_\tri(\Tot(F))\leq 2\nu_\tri(\Tot(F)).
\]
\end{theorem}

\begin{corollary}\label{cor:cubic}
If $F$ is a finite simple cubic triangle-free graph, then
\[
  \tau_\tri(\Tot(F))=\nu_\tri(\Tot(F))=\frac52|V(F)|.
\]
\end{corollary}

The two proofs share an edge decomposition.  Incidences at a root vertex
induce a complete graph, and these complete graphs have pairwise disjoint
edge sets.  Elementary complete-graph packings give local lower bounds.
A single coloring of the root edges controls triangles that cross several
incidence cliques.  In total graphs we also select all original root edges;
the corresponding bridge triangles supply a compatible packing contribution.
A Hall assignment gives a sharp cover when the root is cubic and triangle-free.

The line-graph inequality for triangle-free roots appears in
\cite[Corollary~3.3.3]{munaro2016}; Theorem~\ref{thm:line} removes that
hypothesis.  Nearby results establish Tuza's inequality under a maximum
average degree bound \cite{puleo2015}; a recent preprint treats maximum
degree at most seven \cite{gupta2026}.  Other results improve the coefficient
for certain $K_4$-free graphs \cite{munaro2018}.  Those hypotheses do not cover
the unbounded-degree classes considered here.  We are not aware of a previous
statement of either theorem in the full scope given above.

All graphs in the paper are finite and simple; empty and disconnected roots
are allowed.  Simplicity ensures that two distinct root edges have at most
one common endpoint.  The proofs construct an actual transversal $C$ and
packing $\mathcal P$ in the same graph with $|C|\leq2|\mathcal P|$.
The complete-graph bound uses a three-part construction and strong induction,
with explicit certificates for orders below 21.  These finite certificates
are checked by Lean's kernel; no exact formula for complete-graph packing
numbers is assumed.  The results concern the two stated graph classes and
do not resolve Tuza's conjecture for arbitrary graphs.

\section{Complete-graph tools}\label{sec:clique}

For $d\geq 0$, put
\[
  t(d)=\binom d2-\left\lfloor\frac{d^2}{4}\right\rfloor,
  \qquad D(d)=\nu_\tri(K_d).
\]
A balanced bipartition of $V(K_d)$, followed by the selection of every edge
inside a part, gives a triangle transversal of size $t(d)$.  Explicitly,
\begin{equation}\label{eq:t-values}
  t(2r)=r(r-1),\qquad t(2r+1)=r^2.
\end{equation}

We use the following elementary lower bounds on $D(d)$.  The basic
complete-graph inequality is classical; compare the construction attributed
to Seb\H{o} in \cite[Theorem~3.3.2]{munaro2016}.  The following proof uses the
same Latin-square recurrence and finite bases as the Lean development and
supplies the additional even-order slack.

\begin{lemma}[Local clique slack]\label{lem:clique}
For every integer $d\geq 0$,
\begin{equation}\label{eq:basic-slack}
  t(d)\leq 2D(d).
\end{equation}
If $d\geq 6$ is even, then
\begin{equation}\label{eq:even-slack}
  t(d)+2\leq 2D(d).
\end{equation}
\end{lemma}

\begin{proof}
For $q>0$, partition $3q$ vertices into three copies $A,B,C$ of
$\mathbb Z/q\mathbb Z$.  Take the $q^2$ triangles
\[
  \{a_A,b_B,(a+b)_C\},\qquad a,b\in\mathbb Z/q\mathbb Z.
\]
Any two coordinates determine the third, so distinct triangles share at
most one vertex.  They use only edges between parts.  Adding a maximum
packing inside each of $A,B,C$ therefore gives
\begin{equation}\label{eq:latin-recurrence}
  D(3q)\geq q^2+3D(q).
\end{equation}
Also, $D(a)\leq D(b)$ for $a\leq b$, by embedding $K_a$ into $K_b$.

For completeness, the following describes precisely the finite witnesses
used for $0\leq n<21$, on vertices $0,\ldots,n-1$.  Use the empty list
for $n\leq2$.  For $n\geq3$, put $r=\lfloor n/2\rfloor$ and begin with
\[
  \bigl\{\{i,j,r+((i+j)\bmod r)\}:0\leq i<j<r\bigr\}.
\]
If $n=2r\geq6$, add $\{r,r+1,r+2\}$.  If $n=2r+1$, add
\[
  \{2r,i,r+((2i)\bmod r)\}\quad (i\in I_r),\qquad
  I_r=\begin{cases}
    \{0,\ldots,r-1\},&r\text{ odd},\\
    \{0,\ldots,r/2-1\},&r\text{ even}.
  \end{cases}
\]
Each listed set has three distinct vertices.  For the initial triangles,
sharing a vertex in each of the two halves forces the remaining index to
agree.  The extra even-order triangle uses only edges in the second half.
For odd orders, multiplication by two is injective on $I_r$ modulo $r$;
moreover, an initial triangle containing $i$ cannot also contain
$r+((2i)\bmod r)$, since its other first-half index would then equal $i$.
Thus every pair of distinct listed triangles shares at most one vertex.
Their sizes $p_n$ are
\[
\begin{array}{c|rrrrrrrrrrr}
n&0&1&2&3&4&5&6&7&8&9&10\\\hline
p_n&0&0&0&1&1&2&4&6&7&8&11
\end{array}
\]
and
\[
\begin{array}{c|rrrrrrrrrr}
n&11&12&13&14&15&16&17&18&19&20\\\hline
p_n&15&16&18&22&28&29&32&37&45&46
\end{array}
\]
Direct substitution verifies $t(n)\leq2p_n$ and, for even $n\geq6$,
$t(n)+2\leq2p_n$.  The formal proof checks the same literal lists for
three-element cardinality, pairwise intersection size at most one, and
these numerical bounds.

Proceed by strong induction.  For $n\geq21$, write $n=3q+s$ with
$q\geq7$ and $0\leq s<3$.  The function $t$ is nondecreasing, and
\begin{equation}\label{eq:latin-arithmetic}
  t(3q+s)+2\leq2q^2+3t(q).
\end{equation}
Indeed it suffices to take $s=2$.  When $q=2r$, the right side minus
the left side is $2r^2-6r-2\geq0$ for $r\geq4$; when $q=2r+1$, it
is $2r^2-4r-4\geq0$ for $r\geq3$.  The induction hypothesis,
\eqref{eq:latin-recurrence}, and monotonicity now give
\[
  t(n)+2\leq2q^2+3t(q)
  \leq2\bigl(q^2+3D(q)\bigr)\leq2D(3q)\leq2D(n).
\]
This proves both assertions, including the required even-order slack.
\end{proof}

\section{Line graphs}\label{sec:line}

Let $H$ be a finite simple graph.  For $v\in V(H)$, let $W_v$ be the set of
edges of $L(H)$ corresponding to unordered pairs of distinct root edges
incident with $v$.  Two adjacent edges of a simple graph have a unique common
endpoint, and hence
\begin{equation}\label{eq:wedge-partition}
  E(L(H))=\mathop{\dd}_{v\in V(H)}W_v.
\end{equation}
We call $W_v$ the \emph{wedge packet} at $v$.

\newpage
\begin{lemma}[Line-graph triangle classification]\label{lem:line-class}
Every triangle of $L(H)$ is exactly one of the following:
\begin{enumerate}[label=\textup{(\roman*)}]
  \item a star triangle formed by three root edges incident with one vertex
        $v$, whose three edges all belong to $W_v$;
  \item a root triangle formed by the three edges of a triangle $xyz$ in
        $H$, with one edge in each of $W_x,W_y,W_z$.
\end{enumerate}
\end{lemma}

\begin{proof}
Three distinct pairwise intersecting two-element sets either have a common
element or are $\{x,y\},\{x,z\},\{y,z\}$.  Applied to three pairwise adjacent
root edges, these are precisely the two stated cases.
\end{proof}

The star triangles in $W_v$ form the triangle system of $K_{d_H(v)}$.
Choose $D(d_H(v))$ edge-disjoint star triangles at every root vertex.  The
packet partition \eqref{eq:wedge-partition} shows that their union is a
packing of size
\begin{equation}\label{eq:line-packing}
  P_0=\sum_{v\in V(H)}D(d_H(v)).
\end{equation}
Write
\begin{equation}\label{eq:line-T}
  T=\sum_{v\in V(H)}t(d_H(v)).
\end{equation}
Lemma~\ref{lem:clique} gives $T\leq 2P_0$.

The cover uses a balanced red-blue coloring of the root edges.  We isolate
the required form.

\begin{lemma}[Alternating Euler coloring]\label{lem:euler-color}
Let $K$ be a connected graph.
\begin{enumerate}[label=\textup{(\alph*)}]
  \item If $K$ has an odd-degree vertex, its edges can be colored red and blue
        so that the two color degrees are equal at every even-degree vertex
        and differ by one at every odd-degree vertex.
  \item If every degree of $K$ is even and $|E(K)|$ is even, the two color
        degrees can be made equal at every vertex.
  \item If every degree of $K$ is even and $|E(K)|$ is odd, the two color
        degrees can be made equal except at one chosen seam vertex, where they
        differ by two.
\end{enumerate}
\end{lemma}

\begin{proof}
For (a), add a new dummy vertex adjacent to every odd-degree vertex.  The
augmented graph is connected and Eulerian.  Start an Euler tour at the dummy
vertex and color successive tour edges alternately.  At every original vertex
the tour edges pair into differently colored entering and leaving edges.  Any
odd cyclic seam occurs at the dummy vertex.  After deleting the dummy edges,
the stated balances remain.

For (b) and (c), an isolated vertex needs no coloring.  Otherwise alternately
color an Euler tour of $K$.  If the tour length is even, every visit pairs two
different colors.  If it is odd, the first and
last edges have the same color and the unique discrepancy of two occurs at
the starting vertex.
\end{proof}

Given such a coloring, select in $W_v$ every wedge whose two root edges have
the same color.  Every star triangle is hit by the pigeonhole principle.
Every root triangle is hit because two of its three root edges have the
same color.  If the red and blue degrees at $v$ are $a$ and $b$, the local
cost is
\[
  \binom a2+\binom b2.
\]
By \eqref{eq:t-values}, this equals $t(d_H(v))$ when an even degree is
balanced or an odd degree differs by one.  At a difference-two even degree it
equals $t(d_H(v))+1$.

Consequently, every connected component has a cover of cost at most its
contribution to $T$, except possibly a connected Eulerian component with an
odd number of edges.  Such a component has a cover with one extra edge.  The
following lemma absorbs this unique surcharge.

\begin{lemma}[Odd Eulerian seam repair]\label{lem:seam}
Let $K$ be a connected Eulerian component of $H$ with an odd number of edges.
Put
\[
  T_K=\sum_{v\in V(K)}t(d_H(v)),\qquad
  P_K=\sum_{v\in V(K)}D(d_H(v)).
\]
There is a triangle transversal $C_K$ of $L(K)$ and a triangle packing
$\mathcal P_K$ in $L(K)$ such that $|C_K|\leq 2|\mathcal P_K|$.
\end{lemma}

\begin{proof}
The alternating Euler cover has size at most $T_K+1$.  We distinguish three
exhaustive cases.

If some vertex has degree at least six, its degree is even and
\eqref{eq:even-slack} supplies two units of local packing slack.  Together
with \eqref{eq:basic-slack} at every other vertex, this gives
$T_K+1\leq 2P_K$.

Suppose next that $K$ is triangle-free.  Discard the global coloring.  At
each vertex, split its incident root edges into two equal parts and select the
wedges lying inside a part.  Since there are no root triangles, these
local selections cover every triangle and have total size exactly $T_K$.
Thus $T_K\leq 2P_K$.

It remains that every positive degree is two or four and that $K$ contains a
root triangle $xyz$.  Keep the alternating Euler cover of size at most
$T_K+1$.  At a degree-two vertex the local packing is empty.  At a
degree-four vertex, $D(4)=1$, and a triangle of the local $K_4$ can avoid any
one prescribed wedge: if that wedge is $ab$, take the triangle on
$\{a,c,d\}$.  We may therefore choose the local maximum packings at
$x,y,z$ to avoid, independently, the three wedges of the triangle in $L(K)$
corresponding to $xyz$.  Adding it to the local packing gives a packing of size
$P_K+1$.  Hence
\[
  T_K+1\leq 2P_K+1\leq 2(P_K+1).
\]

These cases are exhaustive because, in the absence of a degree at least six,
Eulerianity forces every positive degree to be two or four, after which the
component is either triangle-free or contains a root triangle.
\end{proof}

\begin{proof}[Proof of Theorem~\ref{thm:line}]
Apply the alternating Euler cover to every component not requiring the seam
repair, and Lemma~\ref{lem:seam} to every connected Eulerian odd-edge
component.  Line graphs and their triangle systems are disjoint across root
components, so the component covers and packings may be united.  This proves
$\tau_\tri(L(H))\leq2\nu_\tri(L(H))$.

For sharpness, $L(K_{1,4})\cong K_4$ has
$(\nu_\tri,\tau_\tri)=(1,2)$.  Indeed, every two of its triangles share an
edge, one edge cannot meet all four triangles, and two opposite edges do.
\end{proof}

\begin{remark}
The quantifier in Lemma~\ref{lem:seam} is one surcharge per connected
Eulerian odd-edge component, not per block, root triangle, or high-degree
vertex.  Odd cycles show why the triangle-free branch must discard the Euler
seam rather than charge it to a nonexistent local packing.
\end{remark}

\section{Total graphs}\label{sec:total}

Let $F$ be a finite simple graph.  Write
\[
  d_v=d_F(v),\qquad m=|E(F)|.
\]
For every $v\in V(F)$, put
\[
  Q_v=\{v\}\cup\delta_F(v),
\]
where $\delta_F(v)$ is the set of root edges incident with $v$, and an edge
in this set is regarded as its edge-vertex in the total
graph.  The set $Q_v$ induces $K_{d_v+1}$.  Its edges are the incidence
edges from $v$ to its incident edge-vertices and the wedge edges between
pairs of root edges meeting at $v$.  These packet edge sets are pairwise
disjoint.  Original root edges lie outside all of them.
For a root edge $uv$, let $B_{uv}=\{u,uv,v\}$ be its \emph{bridge triangle};
here $uv$ inside the vertex set denotes the edge-vertex.

\begin{proof}[Proof of Theorem~\ref{thm:total}]
Apply Lemma~\ref{lem:euler-color} componentwise to color the root edges red
and blue.  At every odd-degree vertex the color degrees differ by one, and
at every even-degree vertex they differ by at most two.

In each $Q_v$, retain these colors on the edge-vertices and color the original
vertex $v$ with the less frequent color, breaking ties arbitrarily.  Adding
one vertex to the minority side changes a discrepancy of zero, one, or two
into one, zero, or one, respectively.  Thus $Q_v$ is balanced and its
within-color edges number exactly $t(d_v+1)$.

Let $C$ consist of these within-color edges in every $Q_v$, together with
all $m$ original root edges.  This gives a single red-blue coloring of all
vertices of $\Tot(F)$: each edge-vertex keeps its root-edge color, and
each original vertex has its chosen minority color.  Every monochromatic
edge of $\Tot(F)$ belongs to $C$.  Indeed, an original-original edge was
selected outright, an incidence edge belongs to its original vertex's
packet, and a wedge belongs to the packet at the unique common root
endpoint.  Every triangle has a monochromatic edge by the pigeonhole
principle, so $C$ is a triangle transversal.  Counting the packet
contributions and original edges gives
\begin{equation}\label{eq:all-total-cover}
  |C|\leq m+\sum_v t(d_v+1).
\end{equation}

For a packing, take all $m$ bridge triangles and, at each $v$, a packing of
$D(d_v)$ triangles among the edge-vertices in $\delta_F(v)$.  Distinct bridges
use distinct original and incidence edges.  The additional triangles use
only wedge edges, whose packets are mutually disjoint.  Thus their union is
one packing $\mathcal P$ of size
\begin{equation}\label{eq:total-packing}
  |\mathcal P|=m+\sum_vD(d_v)\leq\nu_\tri(\Tot(F)).
\end{equation}
The identity and degree-sum bound
\begin{equation}\label{eq:parity-identities}
  t(d+1)-t(d)=\left\lfloor\frac d2\right\rfloor,
  \qquad
  \sum_v\left\lfloor\frac{d_v}{2}\right\rfloor\leq m
\end{equation}
follow from \eqref{eq:t-values} and the degree sum.  Therefore
\[
  \tau_\tri(\Tot(F))\leq |C|
  \leq 2m+\sum_vt(d_v)
  \leq 2|\mathcal P|
  \leq 2\nu_\tri(\Tot(F)),
\]
where the penultimate inequality is Lemma~\ref{lem:clique}.
\end{proof}

\section{Cubic triangle-free roots}\label{sec:triangle-free}

For a cubic triangle-free root, a distinct choice of one incident edge at
each vertex gives a smaller cover.  We first prove the assignment used in
the formal development.

\begin{lemma}[Distinct incident-edge representatives]\label{lem:hall}
If a finite simple graph $F$ has minimum degree at least two, there is an
injective map $p:V(F)\longrightarrow E(F)$ such that $p(v)$ is incident
with $v$ for every vertex $v$.
\end{lemma}

\begin{proof}
Apply Hall's theorem to the vertex-edge incidence bipartite graph.  For
$S\subseteq V(F)$, let $N(S)$ consist of the root edges incident with at
least one vertex of $S$.  Counting incidences gives
\[
  2|S|\leq\sum_{v\in S}d_F(v)\leq2|N(S)|,
\]
because a root edge has at most two endpoints in $S$.  Hence
$|S|\leq|N(S)|$ for every $S$, which is precisely Hall's condition.
\end{proof}

\begin{proof}[Proof of Corollary~\ref{cor:cubic}]
Put $n=|V(F)|$.  Then $m=3n/2$ and $D(3)=1$, so the packing in
\eqref{eq:total-packing} has $3n/2+n=5n/2$ triangles.

Choose $p$ by Lemma~\ref{lem:hall}.  At each $v$, write its three incident
edges as $p(v),a_v,b_v$.  In the packet $Q_v\cong K_4$, select the
incidence edge $v\,p(v)$ and the wedge $a_vb_v$.  These two edges form a
perfect matching of $Q_v$, so every triangle in the packet meets one of
them.  Also select every original edge outside $p(V(F))$.

For a bridge $B_e$, if $e=p(v)$ for some $v$, its selected incidence edge
$v\,p(v)$ meets the bridge; otherwise its original edge is selected.
Since $F$ is triangle-free, every total triangle is a packet triangle or a
bridge.  To see this directly, three original vertices would give a root
triangle; two original vertices force their joining edge-vertex; one
original vertex forces the other two edge-vertices into its packet; and
three edge-vertices are classified by Lemma~\ref{lem:line-class}, whose
root-triangle case is excluded.  This is the same case analysis used in
the formal cover proof.  The selection is therefore a transversal.  Its
size is at most
\[
  2n+\bigl(m-|p(V(F))|\bigr)=2n+(m-n)=m+n=5n/2,
\]
using the injectivity of $p$.  Finally, every transversal contains a
distinct edge for each member of an edge-disjoint packing, so
\[
  \frac{5n}{2}\leq\nu_\tri(\Tot(F))
  \leq\tau_\tri(\Tot(F))\leq\frac{5n}{2}.\qedhere
\]
\end{proof}

\section{Formal verification}\label{sec:formal}

The accompanying \href{https://github.com/samuelmasse/tuza-line-total-graphs}
{Lean repository} implements the constructions presented above.  The local
clique proof uses the Latin packing, the same finite lists for orders below
21, and the same strong-induction inequality.  The line and total proofs
use the Euler colorings, packet unions and seam repair described here;
the cubic cover uses the incident-edge assignment of Lemma~\ref{lem:hall}.
The earlier balanced-orientation refinement is retained in a separate
informal note in the repository and is not a result of this revised paper.

The four headline declarations are
\begin{center}
\small
\begin{tabular}{ll}
Theorem~\ref{thm:line} & \texttt{Tuza.lineGraph\_satisfiesTuza}\\
Line-graph sharpness & \texttt{Tuza.lineGraph\_factor\_sharp}\\
Theorem~\ref{thm:total} & \texttt{Tuza.totalGraph\_satisfiesTuza}\\
Corollary~\ref{cor:cubic} & \texttt{Tuza.cubic\_triangleFree\_total\_parameters}
\end{tabular}
\end{center}
The formal graph definitions use ordinary finite simple graphs, edge-disjoint
triangle packings and edge transversals.  Empty and disconnected roots are
included.  The cubic conclusion uses the equivalent natural-number identities
$2\tau_\tri=5|V(F)|$ and $2\nu_\tri=5|V(F)|$.

At source commit \texttt{82f6bb5}, using Lean 4.33.1 and the pinned Mathlib
dependency, all four declarations passed
\href{https://github.com/PalomarRegistry/PalomarSubmission/actions/runs/34193414990}
{Palomar's Comparator check and replay by Lean and NanoDa}.
Only propositional extensionality, classical choice and quotient soundness
occur in the audited axioms.  The deliberately incomplete Challenge
statements are separate from the completed Solution proofs.  A detailed
proof-to-declaration map and full commit identifiers are provided in the
repository.  That mechanical run checked the Lean sources before this
expository revision; it did not review this revised manuscript, establish
novelty, or constitute human expert review or completed registration.

\section{Concluding remarks}

The two operator classes permit globally compatible certificates because
adjacent root edges have a unique common endpoint.  In line graphs, the
only possible loss in an Euler coloring is one extra selected wedge per
connected Eulerian component with an odd number of edges.  In total graphs,
the original vertex balances its enlarged incidence clique, while bridge
triangles pay for the selection of original root edges.

The factor two is sharp for line graphs.  We do not determine the best
factor for total graphs, although the cubic triangle-free case has equality
between packing and transversal numbers.  Extensions to multigraph roots,
generalized line graphs, or quasi-line graphs require additional arguments
and are not claimed here.

\paragraph{AI assistance.}
The author acknowledges the assistance of OpenAI's GPT-6 Astra in developing
the proofs, checking arguments, searching the literature, preparing the
exposition, and constructing the Lean formalization.  The finite clique
certificates are part of the proof and are checked by Lean's kernel;
exploratory computations and bounded graph tests are not premises of the
universal results.

\bibliographystyle{plain}
\begingroup
\raggedright
\bibliography{references}
\endgroup

\end{document}
